\IfFontExistsTF{Noto Serif CJK SC}{
  \setCJKmainfont{Noto Serif CJK SC}[BoldFont={Noto Serif CJK SC Bold},AutoFakeSlant=0.15]
  \setCJKsansfont{Noto Sans CJK SC}
  \setCJKmonofont{Noto Sans Mono CJK SC}
}{
  \setCJKmainfont{FandolSong-Regular.otf}[BoldFont=FandolSong-Bold.otf,ItalicFont=FandolKai-Regular.otf]
  \setCJKsansfont{FandolHei-Regular.otf}[BoldFont=FandolHei-Bold.otf]
  \setCJKmonofont{FandolFang-Regular.otf}
}
\usepackage[a4paper,margin=25mm,headheight=15pt]{geometry}
\usepackage{amsmath,amssymb,amsthm,bm,mathtools}
\usepackage{float}
\usepackage{graphicx,booktabs,longtable,array,tabularx,multirow}
\usepackage{enumitem,xcolor,tikz,pgfplots}
\usetikzlibrary{arrows.meta,positioning,calc,shapes.geometric}
\pgfplotsset{compat=1.18}
\usepackage{fancyhdr}
\usepackage[unicode,colorlinks=true,linkcolor=blue!45!black,citecolor=green!35!black,urlcolor=blue!55!black]{hyperref}
\usepackage{bookmark}
\usepackage{listings}
\lstset{basicstyle=\ttfamily\footnotesize,breaklines=true,frame=single,columns=fullflexible,keepspaces=true,showstringspaces=false,language=Python}
\setlist{nosep,topsep=5pt,leftmargin=2em}
\setlength{\parskip}{3pt}
\setlength{\emergencystretch}{2em}
\renewcommand{\arraystretch}{1.18}
\newtheorem{theorem}{定理}[chapter]
\newtheorem{proposition}[theorem]{命题}
\newtheorem{lemma}[theorem]{引理}
\theoremstyle{definition}
\newtheorem{definition}[theorem]{定义}
\newtheorem{example}[theorem]{例}
\newtheorem{remark}[theorem]{说明}
\newcommand{\R}{\mathbb R}
\newcommand{\C}{\mathbb C}
\newcommand{\E}{\mathbb E}
\newcommand{\Prob}{\mathbb P}
\newcommand{\HH}{\mathsf H}
\newcommand{\T}{\mathsf T}
\newcommand{\I}{\mathbf I}
\newcommand{\one}{\mathbf 1}
\newcommand{\norm}[1]{\left\lVert #1\right\rVert}
\newcommand{\abs}[1]{\left\lvert #1\right\rvert}
\newcommand{\pos}[1]{\left[#1\right]_+}
\DeclareMathOperator{\tr}{tr}
\DeclareMathOperator{\diag}{diag}
\DeclareMathOperator{\rank}{rank}
\DeclareMathOperator{\Rea}{Re}
\DeclareMathOperator{\Ima}{Im}
\DeclareMathOperator{\SINR}{SINR}
\DeclareMathOperator{\prox}{prox}
\DeclareMathOperator{\dist}{dist}
\DeclareMathOperator*{\argmin}{arg\,min}
\DeclareMathOperator*{\argmax}{arg\,max}
\newcommand{\chaptertakeaway}[1]{\begin{quote}\small\textbf{本章目标}\quad #1\end{quote}}
\pagestyle{fancy}
\fancyhf{}
\fancyhead[L]{\small 通信中的优化问题}
\fancyhead[R]{\small\nouppercase{\leftmark}}
\fancyfoot[C]{\thepage}
\ctexset{chapter={format=\huge\bfseries,beforeskip=8pt,afterskip=22pt},section={format=\Large\bfseries}}
\numberwithin{equation}{chapter}
\graphicspath{{figures/}}
\hypersetup{pdftitle={通信中的优化问题},pdfsubject={通信优化建模、算法与可复现实验},pdfauthor={}}
